\SourceLeaf{384}%
drate war viel geringer, als sie bei regelmässig vorausgesetzter Erdfigur sein sollte. Die neue Messung hat gezeigt, dass \textsc{Beccaria} bei der astronomischen Bestimmung dieser Krümmung allerdings einen Fehler von \(13''41\) begangen hat (der bei der Unvollkommenheit seiner Instrumente sehr verzeihlich ist); allein das Zeichen dieses Fehlers ist das entgegengesetzte von dem vermutheten, und die Anomalie wird also noch um so viel vergrössert. Die nach obigen Elementen aus den geodaetischen Messungen berechnete Amplitudo ist nemlich (auf \textsc{Beccarias} Endpunkte reducirt) \(1^0\,8'\,18''91\); die aus \textsc{Beccarias} astronomischen Beobachtungen sich ergebende \(1^0\,7'\,44''30\), und die neue Bestimmung \(1^0\,7'\,31''07\). Die neue Messung ist mit so guten Hülfsmitteln und mit so ausgezeichneter Sorgfalt ausgeführt, dass man gezwungen ist, diesen grossen Unterschied von \(47''84\) fast ganz als eine Unregelmässigkeit der Erdfigur zu betrachten, die merkwürdigste Thatsache dieser Art, die bisher in den Annalen der höhern Geodaesie vorgekommen ist. Höchst wahrscheinlich ist die Attraction der diese Messung in Norden und Süden begrenzenden Alpenketten eine Hauptursache dieses Phänomens, allein eben so wahrscheinlich hat die ungleiche Dichtigkeit der untern Erdschichten, vielleicht bis zu grosser Tiefe hinab, nicht minder Antheil daran. Wenigstens lassen sich ähnliche bei ganz in der Ebene liegenden Punkten vorgekommene Unterschiede von sehr bedeutender Grösse (z.\,B. eine Anomalie von \(21''9\) zwischen Mailand und Parma) nicht wohl anders erklären. Wir setzen hinzu, dass je mehr die sorgfältig ausgeführten Gradmessungen vervielfältigt werden, desto mehr die Ueberzeugung Platz gewinnt, dass solche Abweichungen nur in Rücksicht auf ihre Grösse, aber nicht an sich als Ausnahmen betrachtet werden dürfen, und dass sich solche nach grösserm oder kleinerm Massstabe überall zeigen. Die Verf. haben eine interessante vergleichende Übersicht der durch astronomische Beobachtungen bestimmten und der durch geodaetische Messungen berechneten Polhöhen von 34 über halb Europa zerstreuten und durch Dreiecke unter sich verbundenen Punkten gegeben. Freilich hat man dieselbe nur wie einen unvollkommenen Versuch zu betrachten, da sie grossentheils nur auf unbeglaubigten fragmentarischen Notizen von den Resultaten der geodaetischen Messungen beruht: denn leider sind die meisten dieser Messungen in Frankreich, Italien, Oesterreich und Baiern noch immer nicht bekannt gemacht.

Derselbe Abschnitt enthält ausserdem noch die neue Messung einer kleinen Grundlinie bei Turin, wodurch einige Umstände, welche die von Hrn.\ \textsc{von}
